\chapter*{Abstract}\addcontentsline{toc}{chapter}{Abstract}
This project focuses on trajectory, measurement, contact, and support issues in robotic motion control, progressively completing ROS2 simulation communication, Panda joint trajectory tracking, inertial measurement unit (IMU) filtering and physical data acquisition, end-effector admittance control, inverted pendulum balancing, planar bipedal walking, and three-dimensional humanoid whole-body control. The technical approach begins with simplified scenarios that isolate individual problems, gradually reintroducing gravity, collisions, floating base, and limited motor torque, while using failure cases to validate the applicability scope of each solution.

For the manipulator, drive-gain tuning and reference velocity reduced mean motion error on independent-amplitude tests from 0.864 degrees to 0.082 degrees for a 5-second quintic trajectory. The physical IMU on Jetson produced new measurements at about 25Hz; at 200Hz publication, about 87.52\% of messages reused earlier measurements. Publication frequency and information update rate must therefore be assessed separately. Contact control progressed from improving mean force to approach deceleration, target-force ramping, and gravity compensation. Under the selected rigid-contact condition at a 1ms step, the raw peak was 11.11N and final mean force about 5.007N. Halving the physics step raised the raw peak to 20.43N, showing remaining numerical sensitivity.

For bipedal control, constrained foot placement and causal delay prediction were first tested in a two-dimensional linear inverted pendulum model. Continuous swing-leg dynamics, double support, and inverse-dynamics quadratic programming were then connected to articulated MuJoCo dynamics. The later planar model completed two twenty-step walks, each advancing about 2.323m; 79 of 94 unique conditions passed, with frequent failures at longer delays. The three-dimensional G1 retained a free-floating base and 29 torque-limited motors for standing, single-foot support on both sides, and slow alternating steps. Of 52 formal conditions, 51 passed; only 2 of 8 additional strong-disturbance conditions passed with a frozen controller. The four-step reference took 29 seconds and advanced about 0.210m, representing quasistatic walking validation.

The report is organized into seven technical sections and twenty-eight stages, preserving the model assumptions, metric definitions, parameter comparisons, failure boundaries, and code entry points for each stage. Hardware measurement results are limited to explicitly annotated IMU and Jetson data collection; mechanical arm contact, balance, and walking results are all from simulation. The current stage establishes a control method implementation and physical validation pipeline, which does not equate to real robot deployment or hard real-time control systems.

\noindent\textbf{Keywords:} Robot control; State feedback; Admittance control; Inverted pendulum; Whole-body inverse dynamics; Bipedal walking; Physical simulation.

\paragraph{This Supplement} After the original 28 stages, we add the official G1 reinforcement learning walking reproduction and verification of contact, direction, and stopping, totaling 29 stages. We add 12 fixed scenarios amounting to 245 seconds of simulation and 13 video clips. The official policy was neither trained nor modified; it uses a different model than the self-developed quasistatic controller and is not subject to fair algorithm ranking. Dynamic stepping has been reproduced, but heading offset and continuous stepping after zeroing still require separate handling.
